\documentclass[11pt,a4paper]{article}

% --- Sprache & Encoding ---
\usepackage[ngerman]{babel}
\usepackage[T1]{fontenc}
\usepackage[utf8]{inputenc}

% --- Layout ---
\usepackage[margin=2.5cm]{geometry}
\usepackage{setspace}
\onehalfspacing

% --- Tabellen & Listen ---
\usepackage{booktabs}
\usepackage{longtable}
\usepackage{enumitem}
\setlist[itemize]{leftmargin=1.5em, itemsep=0.2em, topsep=0.3em}

% --- Code ---
\usepackage{listings}
\usepackage{xcolor}
\lstdefinestyle{code}{
    basicstyle=\ttfamily\small,
    breaklines=true,
    frame=single,
    columns=fullflexible,
    keywordstyle=\color{blue!60!black}\bfseries,
    commentstyle=\color{gray},
    stringstyle=\color{green!40!black},
    numbers=left,
    numberstyle=\tiny\color{gray},
    tabsize=4,
}
\lstset{style=code}

% --- Grafik & Diagramme ---
\usepackage{graphicx}
\usepackage{tikz}
\usetikzlibrary{shapes,arrows.meta,positioning}

% --- Mathe ---
\usepackage{amsmath}
\usepackage{amssymb}

% --- Links (immer als letztes Paket) ---
\usepackage[hidelinks]{hyperref}

\title{Abaqus Handbuch-Chatbot \\ \large Ein Retrieval-Augmented-Generation-System zur Beantwortung von Bedienungsfragen zur FEA-Software Abaqus}
\author{
    Florian Heithoff \\[0.3em]
    \small Matrikelnummer: 30483505 \\
    \small NLP-Studienmodul \\
    \small Fachhochschule Südwestfalen \\
    \small Betreuer: Prof. Dr. C. Gawron
}
\date{03. September 2026}

\begin{document}
% Deaktiviert babel/ngerman's "-Shorthands (z. B. "a, "s, "I) - das Dokument
% nutzt durchgängig direkte UTF-8-Umlaute/ß und verwendet " stattdessen als
% gewöhnliches Anführungszeichen (sonst würde z. B. "Ist ... zu Ïst).
\shorthandoff{"}

\maketitle
\tableofcontents
\newpage

% =====================================================================
\section{Einleitung}
% =====================================================================

\subsection{Über Abaqus}

Abaqus ist eine Simulationssoftware aus der Strukturmechanik, die mithilfe der
Finite-Elemente-Methode (FEM) strukturmechanische Problemstellungen numerisch
löst. Eingesetzt wird sie vor allem in der Konstruktion und Entwicklung von
mechanisch belasteten Bauteilen. Sie erlaubt es, deren Verhalten unter Last bereits am digitalen Modell zu prüfen, bevor ein physischer
Prototyp gebaut wird. Abaqus ist eine komplexe Software mit zahlreichen Modulen, Optionen und einer umfangreichen
Keyword-/Input-Deck-Syntax (siehe \texttt{Abaqus2017\_KEYWORDS.pdf}), deren
sichere Bedienung entsprechend anspruchsvoll ist.

\subsection{Motivation und Zielsetzung}

Die offiziellen Abaqus-Handbücher umfassen mehrere tausend Seiten, verteilt
auf mehrere PDF-Dokumente. Eine reine Volltextsuche im PDF-Reader findet
zwar Fundstellen, liefert aber keine kontextuelle Antwort auf eine konkrete
Bedienungsfrage. Ziel dieser Arbeit ist daher ein Chatbot, der Fragen in
natürlicher Sprache beantwortet, und zwar ausschließlich auf Basis der
offiziellen Handbücher und nicht aus dem "Weltwissen" des zugrunde liegenden
Sprachmodells. Zwei Eigenschaften sind für die Antwortqualität zentral: Eine
seitengenaue Quellenangabe, die die Antwort im Original nachprüfbar macht,
und ein Gesprächsgedächtnis, das Folgefragen versteht. Halluzinationen wird
dabei durch einen strikten System-Prompt entgegengewirkt, sobald der abgerufene Kontext die
Frage nicht abdeckt.

Über diese Grundfunktion hinaus ist jede Komponente der Verarbeitungspipeline
-- Parser, Chunking, Embedding, Retrieval und Antwort-LLM -- bewusst
austauschbar gestaltet. Das erlaubt es, für jede Komponente empirisch zu
prüfen, welche Variante die beste Antwortqualität liefert.

% =====================================================================
\section{Vorgehen}
\label{sec:vorgehen}
% =====================================================================

\subsection{Methodisches Vorgehen}

Das Projekt folgt einem iterativen Vorgehen in drei Phasen. Zunächst wird
eine funktionierende Produktiv-Baseline mit allen fünf Pipeline-Stufen --
Parser, Chunking, Embedding, Retrieval und Antwort-LLM -- aufgebaut. Darauf
aufbauend wird systematisch je eine Pipeline-Dimension gegen diese Baseline
variiert und die Auswirkung quantitativ per RAGAS gemessen. Abschließend
werden die jeweiligen Einzelsieger im Best-of-Breed-Experiment
zusammengeführt. Jede Designentscheidung ist dadurch nicht nur begründet,
sondern nachträglich auch empirisch geprüft.

\subsection{Design-Prinzip: isolierte Ein-Dimensions-Vergleiche}

Jede der sechs Vergleichsstudien variiert genau eine Dimension, während alle
übrigen Pipeline-Stufen auf dem jeweiligen Produktiv-Default fixiert bleiben.
Ziel ist es, gemessene Effekte einzelnen Ursachen zurechenbar zu halten,
statt sie in einer vollen Kombinatorik zu vermischen, da mit jeweils vier bis sieben Varianten eine vollständige Kombinatorik zu aufwändig wäre.
Wechselwirkungen zwischen mehreren gleichzeitig veränderten Dimensionen werden dadurch nicht
erfasst und werden erst im abschließenden Best-of-Breed-Experiment
(Abschnitt~\ref{sec:best-of-breed}) sichtbar.

\subsection{Demo-Korpus vs. Produktivkorpus}

Der Produktivkorpus umfasst vier vollständige Abaqus-Handbücher mit ca.
5.100 Seiten und 15.408 Chunks. Alle sechs Vergleichsstudien laufen jedoch
nicht auf diesem vollen Korpus, sondern auf einem festen, deterministisch
zusammengestellten Demo-Korpus von ca. 53 Seiten. Der kleinere Korpus wurde gewählt, da Unstructured (\texttt{hi\_res}) auf dem vollen
Korpus mit geschätzt mehreren Stunden Laufzeit nicht praktikabel wäre --
Docling ist mit warmem Modell-Cache zwar ebenfalls deutlich langsamer als
PyMuPDF4LLM, mit geschätzt rund einer Stunde für den vollen Korpus aber
grundsätzlich machbar (siehe \texttt{docs/PARSER.md}). Zusätzlich wären auch RAGAS-Läufe mit LLM-Richter über
mehrere Kombinationen hinweg  auf dem vollen Korpus zeit- und
kostenintensiv. Um trotz der reduzierten Größe aussagekräftig zu bleiben,
wurde der Demo-Korpus gezielt aus genau den Inhaltstypen zusammengesetzt,
die sich zuvor als unterschiedlich schwierig erwiesen hatten -- darunter
TOC-Tabellen, linienlose Mini-Tabellen, ein vollständiges Tutorial-Kapitel
sowie ein Analyse-Kapitel.

% =====================================================================
\section{Das RAG-System: Architektur}
\label{sec:architektur}
% =====================================================================

\subsection{Systemüberblick}

Das System besteht aus zwei vollständig entkoppelten Phasen. In der
\textbf{Ingestion} -- einmalig ausgeführt bzw. wiederholt bei geänderten
PDFs -- durchlaufen die Handbücher eine feste Kette: Ein Parser wandelt die
PDFs in Text um, ein Chunker zerlegt diesen Text in handhabbare Abschnitte,
ein Embedding-Modell wandelt jeden Abschnitt in einen Vektor um, und das
Ergebnis landet in einem persistenten Chroma-Vectorstore. In der
\textbf{Chat-Session} greift der Bot ausschließlich lesend auf diesen fertigen
Vectorstore zu -- ohne Rückwirkung auf die Wissensbasis während des Chats.
Eine Nutzerfrage durchläuft dabei zunächst einen Contextualize-Step, der
Folgefragen anhand des Gesprächsverlaufs eigenständig verständlich macht,
wird anschließend vom Retriever gegen den Vectorstore aufgelöst und
zusammen mit dem gefundenen Kontext an eine Answer-Chain übergeben, die aus
System-Prompt, Kontext und Frage die endgültige Antwort samt
Quellenangaben generiert.

\subsection{Zentrale Designentscheidungen}

Neben den fünf austauschbaren Pipeline-Stufen (Abschnitt~\ref{sec:untersuchte-varianten})
enthält das System eine Reihe infrastruktureller Entscheidungen, die
bewusst \emph{nicht} zum Gegenstand eines Vergleichs gemacht wurden, weil
sie den Kern der RAG-Pipeline selbst nicht betreffen, sondern deren
technisches Fundament bilden. Sie sind für alle in
Abschnitt~\ref{sec:vergleiche} untersuchten Varianten gleichermaßen
gültig und werden in keiner der sechs Vergleichsstudien verändert:

\begin{itemize}
    \item \textbf{Vektor-DB}: Chroma (lokal, persistent, SQLite-basiert) -- kein Serverbetrieb nötig, native LangChain-Integration, Metadaten-Filterung
    \item \textbf{RAG-Orchestrierung}: LCEL-Runnables statt \texttt{langchain.chains} (in LangChain~$\geq$1.0 entfernt).
    \item \textbf{Conversation Memory}: einfache Nachrichtenliste innerhalb der CLI-Session
\end{itemize}

\subsection{Untersuchte Pipeline-Varianten}
\label{sec:untersuchte-varianten}

Im Gegensatz dazu wurden genau die fünf Stufen, die den eigentlichen
Verarbeitungsweg von der PDF-Seite bis zur fertigen Antwort ausmachen,
bewusst austauschbar gestaltet: Für jede von ihnen existiert im Code eine
Registry mehrerer gleichwertiger Implementierungen (Abschnitt~\ref{sec:code}),
sodass sich jede Stufe unabhängig von den übrigen gegen Alternativen testen
lässt. Diese fünf Stufen sind zugleich die fünf Dimensionen der
Vergleichsstudien aus Abschnitt~\ref{sec:vergleiche}, wo die jeweils
getesteten Varianten und der eingesetzte Produktiv-Default im Detail
aufgeführt sind. An dieser Stelle nur die Übersicht, wie viele Varianten
je Dimension untersucht wurden:

\begin{itemize}
    \item PDF-Parser -- 4 getestete Varianten
    \item Chunking-Strategie -- 6 getestete Varianten
    \item Embedding-Modell -- 7 getestete Varianten
    \item Retrieval-Strategie -- 4 getestete Varianten
    \item Antwort-LLM -- 3 getestete Varianten
\end{itemize}

% =====================================================================
\section{Implementierung}
\label{sec:code}
% =====================================================================

\subsection{Projektstruktur}

Der Kern der Anwendung ist als eigenständiges Python-Paket
\texttt{src/rag\_manual\_bot/} organisiert und darin in vier funktionale
Bereiche unterteilt. \texttt{scripts/}, \texttt{tests/}, \texttt{app.py}
und \texttt{main.py} sind nicht Teil dieses
Pakets, sondern liegen als eigenständige Verzeichnisse bzw. Dateien direkt
im Projekt-Root. Sie importieren das Paket und nutzen dessen Funktionalitäten.

\begin{lstlisting}[numbers=none, basicstyle=\ttfamily\scriptsize, frame=single, breaklines=true]
Abaqus_chatbot/
|-- src/rag_manual_bot/     <- das Python-Paket
|   |-- ingestion/          Loader, Chunker, Parser-/Chunking-/Embedding-Backends
|   |-- rag/                Vectorstore, Retriever, Retrieval-Backends, LCEL-Kette
|   |-- eval/               RAGAS-Anbindung, Evaluations-Datensatz, Metriken
|   |-- citations.py        Quellen-Links
|   `-- cli/chat.py         Terminal-Chat-Oberflaeche
|
|-- scripts/                CLI-Entry-Points fuer Ingestion, Demo-Korpus-Aufbau, RAGAS-Laeufe
|-- tests/                  51 offline lauffaehige Tests
|-- app.py                  Frontend-Entry-Point
`-- main.py                 Terminal-Entry-Point
\end{lstlisting}

\subsection{Zentrale Code-Bausteine}

Dieser Abschnitt beschreibt, wie die in Abschnitt~\ref{sec:architektur}
konzeptionell vorgestellten Komponenten konkret im Code umgesetzt sind --
mit Fokus auf Implementierungsdetails, Funktionsnamen und technische
Kniffe je Baustein.

\subsubsection{Loader}

Die Kernfunktion \texttt{load\_pdf()} in \texttt{ingestion/loader.py}
kombiniert zwei Bibliotheken: \texttt{pymupdf4llm} liefert den
Markdown-Text pro Seite, \texttt{pymupdf} zusätzlich die im PDF
eingebetteten Page-Labels, deren Auflösung über
\texttt{doc[pdf\_index].get\_label()} erfolgt -- mit Fallback auf die
sequenzielle Position, falls kein Label vorhanden ist. Rückgabetyp ist
eine Liste von \texttt{PageDocument}-Objekten (Text, Quelldatei, gedruckte
Seitenzahl); leere Seiten werden bereits hier herausgefiltert, bevor sie
in den Chunker gelangen.

\subsubsection{Chunker}

Die Kernfunktion \texttt{chunk\_pages()} in \texttt{ingestion/chunker.py}
ist nur ein dünner Wrapper um
\texttt{ingestion/chunking\_backends/header\_recursive\_backend.py} mit
den Produktiv-Settings (\texttt{CHUNK\_SIZE}, \texttt{CHUNK\_OVERLAP}) --
dieselbe Backend-Funktion ist zusätzlich unter zwei weiteren
Größenvarianten sowie drei strukturell anderen Chunking-Strategien
registriert (siehe Abschnitt~\ref{sec:code} zum Registry-Pattern).

\subsubsection{Vectorstore}

Die Kernfunktionen \texttt{build\_vectorstore()} und
\texttt{load\_vectorstore()} in \texttt{rag/vectorstore.py} persistieren
lokal über Chroma (SQLite-basiert). Beide tragen einen optionalen
\texttt{embeddings}-Parameter (eine fertige \texttt{Embeddings}-Instanz
statt eines Modellnamens), da der Vectorstore zwingend mit demselben
Embeddings-Objekt geladen werden muss, mit dem er aufgebaut wurde;
\texttt{build\_vectorstore()} löscht vor dem Neuaufbau alle bestehenden
Einträge (\texttt{store.delete(ids=existing\_ids)}) und fügt Dokumente in
Batches von 200 ein.

\subsubsection{Prompts}

\texttt{rag/prompts.py} definiert zwei getrennte Prompt-Templates für die
zwei Aufgaben der Kette: Der Contextualize-Prompt formuliert Folgefragen
anhand des Verlaufs eigenständig um, ohne sie zu beantworten; der
Answer-Prompt legt die Systemrolle fest ("technischer Assistent für
Abaqus"), erzwingt Kontext-Bindung und ein explizites "Ich weiß es nicht"
bei fehlender Abdeckung. Beide binden den bisherigen Gesprächsverlauf über
\texttt{MessagesPlaceholder("chat\_history")} als vollwertige
Chat-Nachrichten ein, nicht als reinen Textblock.

\subsubsection{LCEL-Kette}

Die Kette wird in \texttt{rag/chain.py}, Funktion
\texttt{build\_rag\_chain(vectorstore, llm\_provider,
retrieval\_strategy)}, schrittweise über verkettete
\texttt{RunnablePassthrough.assign(...)}-Aufrufe konstruiert; eine
\texttt{RunnableBranch} überspringt dabei den Contextualize-Schritt bei
leerem Gesprächsverlauf und spart so einen unnötigen LLM-Call. Jeder
\texttt{.assign(...)}-Schritt reichert das durchlaufende Dict um ein
weiteres Feld an (\texttt{standalone\_question}, \texttt{source\_documents},
\texttt{context}, \texttt{answer}). \texttt{llm\_provider} und
\texttt{retrieval\_strategy} werden als Parameter durchgereicht und wählen
aus den jeweiligen Registries (Abschnitt~\ref{sec:code}) -- derselbe
Kettenaufbau bedient damit sowohl den Produktivbetrieb als auch alle
Vergleichsstudien.

\subsubsection{Source-Tracking}

\texttt{format\_docs()} baut aus den abgerufenen Chunks den nummerierten
Kontextblock für das LLM, inklusive Quelle und Seitenzahl als Zitat;
\texttt{citations.py} erzeugt daraus eine klickbare \texttt{file://}-URL
mit PDF-Seitenanker (\texttt{\#page=N}). Der Anker muss sich auf die
physische Seitenposition beziehen statt auf das gedruckte Label, da
PDF-Viewer \texttt{\#page=N} als reinen Dateiindex interpretieren -- die
Pipeline führt daher beide Werte parallel mit. \texttt{\_format\_sources()}
(in \texttt{cli/chat.py} bzw. \texttt{app.py}) dedupliziert die Quellen
abschließend nach Datei und Seite, wobei die Reihenfolge der ersten
Nennung erhalten bleibt.

\subsection{Austauschbare Komponenten (Registry-Pattern)}

Ein einheitliches Implementierungsprinzip zieht sich durch alle fünf
austauschbaren Pipeline-Stufen: Für jede Stufe existiert eine Registry --
ein Dictionary, das einen Namen auf eine Implementierung oder Factory
abbildet. Neue Varianten werden als zusätzlicher Registry-Eintrag ergänzt,
ohne den übrigen Code zu verändern; genau das ermöglicht denselben
Code-Pfad sowohl für den Aufbau der Vergleichs-Collections als auch für
die automatisierten RAGAS-Vergleiche. Konkret:

\begin{itemize}
    \item \texttt{parser\_backends/}: 4 Parser mit identischer Schnittstelle \texttt{parse(pdf\_path) -> dict[int, str]}
    \item \texttt{chunking\_backends/}: 6 Chunking-Strategien
    \item \texttt{embedding\_models.py::EMBEDDING\_BACKENDS}: 7 Embedding-Modelle (3 OpenAI-Factories, 4 lokale HuggingFace-Factories mit Lazy-Import)
    \item \texttt{retrieval\_backends.py::RETRIEVAL\_BACKENDS}: 4 Retrieval-Strategien
    \item \texttt{rag/llms.py::get\_llm(provider)}: 3 Antwort-LLM-Provider (OpenAI, Mistral, GPT-4o)
\end{itemize}

\subsection{Tests}

Die Testsuite (51 Tests) läuft vollständig offline -- es wird kein echter
API-Call gegen OpenAI oder ein anderes Modell ausgeführt. Möglich macht das
\texttt{tests/conftest.py}, das einen Dummy-\texttt{OPENAI\_API\_KEY} setzt,
damit das beim Import validierende \texttt{Settings}-Modul auch ohne
echten Key geladen werden kann, sowie eine selbstgeschriebene, deterministische
\texttt{FakeEmbeddings}-Klasse anstelle echter Embedding-API-Aufrufe.
Abgedeckt sind unter anderem die Chunker-Grundlogik, alle
Chunking-Backends, die Embedding-Modell-Registry, die Retrieval-Backends,
die LLM-Provider-Auswahl, die Custom-Demo-Kombinationslogik sowie der
Eval-Datensatz. Ausführung über \texttt{python -m pytest tests/ -v}.

% =====================================================================
\section{Untersuchte Komponenten}
\label{sec:vergleiche}
% =====================================================================

Für jede der fünf austauschbaren Pipeline-Stufen aus
Abschnitt~\ref{sec:untersuchte-varianten} wurden mehrere konkrete
Implementierungen gegeneinander getestet. Dieses Kapitel listet je
Dimension, \emph{welche} Varianten untersucht wurden und markiert die
jeweilige Produktiv-Default-Wahl; die zugehörige RAGAS-Methodik und die
gemessenen Ergebnisse folgen in Kapitel~\ref{sec:validierung}.

\subsection{PDF-Parser}

Der Parser überführt jede Handbuchseite aus dem PDF in maschinenlesbaren Text.

\begin{itemize}
    \item \textbf{PyMuPDF4LLM (Produktiv-Default)}: wandelt PDF-Seiten seitenweise in Markdown um, erhält dabei Überschriften, Tabellen und Seitenzahlen
    \item \textbf{Docling}: Dokumentenkonvertierungs-Toolkit mit Deep-Learning-basierter Layout- und Strukturerkennung
    \item \textbf{pdfplumber}: Python-Bibliothek mit Fokus auf präzise Extraktion von Text und Tabellen anhand der PDF-Koordinaten
    \item \textbf{Unstructured}: allgemeines Dokument-Partitionierungs-Framework, hier im rechenintensiveren \texttt{hi\_res}-Modus mit eigenem Layout-Erkennungsmodell getestet
\end{itemize}

\subsection{Chunking-Strategien}

Die Chunking-Strategie zerlegt den geparsten Text in einzelne, embedding-fähige Abschnitte.

\begin{itemize}
    \item \textbf{Header+Recursive (Produktiv-Default)}: Hybrid aus Markdown-Header-Split und rekursivem Zeichen-Split als Fallback für überlange Abschnitte; getestet in drei Größenvarianten (400/800/1600 Zeichen)
    \item \textbf{Recursive-only}: rein zeichenbasierte rekursive Aufteilung ohne Berücksichtigung von Überschriften, als naive Baseline
    \item \textbf{Token-basiert}: Chunk-Grenzen anhand der Tokenzahl statt der Zeichenzahl
    \item \textbf{Semantic Chunking}: embedding-basierte Bestimmung der Chunk-Grenzen, setzt Trennstellen dort, wo sich die semantische Ähnlichkeit benachbarter Sätze am stärksten ändert
\end{itemize}

\subsection{Embedding-Modelle}

Das Embedding-Modell überführt jeden Chunk in einen Vektor und bildet damit die Grundlage der semantischen Suche.

\begin{itemize}
    \item \textbf{OpenAI (API-basiert)}: \texttt{text-embedding-3-small} (Produktiv-Default), \texttt{text-embedding-3-large}, \texttt{text-embedding-ada-002}
    \item \textbf{Lokale HuggingFace-/\texttt{sentence-transformers}-Modelle}: \texttt{multilingual-e5-large}, \texttt{bge-m3}, \texttt{paraphrase-multilingual-mpnet}, \texttt{minilm-l6-en} -- laufen lokal, ohne API-Kosten
\end{itemize}

\subsection{Retrieval-Strategien}

Die Retrieval-Strategie bestimmt, welche Chunks zur Anfragezeit aus dem Vectorstore ausgewählt werden.

\begin{itemize}
    \item \textbf{MMR (Produktiv-Default)}: Max Marginal Relevance, wählt aus einem größeren Kandidatenpool eine thematisch diverse Teilmenge
    \item \textbf{Similarity}: reine Nächste-Nachbarn-Suche über Kosinus-Ähnlichkeit, ohne Diversitätsmechanismus
    \item \textbf{Cross-Encoder-Rerank}: sortiert die von der initialen Suche gelieferten Kandidaten mit einem Cross-Encoder-Modell neu
    \item \textbf{BM25-Hybrid}: kombiniert klassische BM25-Stichwortsuche mit dichter Vektorsuche über Reciprocal Rank Fusion
\end{itemize}

\subsection{Antwort-LLMs}

Das Antwort-LLM generiert aus Frage und abgerufenem Kontext die finale Antwort.

\begin{itemize}
    \item \textbf{OpenAI \texttt{gpt-4o-mini} (Produktiv-Default)}
    \item \textbf{Mistral \texttt{mistral-small-latest}}
    \item \textbf{OpenAI \texttt{gpt-4o}}: volle Modellgröße derselben Generation wie der Produktiv-Default, läuft über denselben Hub/Key wie \texttt{gpt-4o-mini}
\end{itemize}

% =====================================================================
\section{Validierung: Quantitative Evaluation mit RAGAS}
\label{sec:validierung}
% =====================================================================

\subsection{Warum RAGAS, mit festem, unabhängigem Richter}

Für diese Arbeit wird daher RAGAS
eingesetzt, ein etabliertes Open-Source-Framework zur automatisierten,
metrikbasierten Evaluation von RAG-Systemen. RAGAS benötigt für jede
Frage vier Bestandteile: die Nutzerfrage selbst, die vom System generierte
Antwort und die dafür abgerufenen Kontext-Chunks -- für drei der fünf in
Abschnitt~\ref{sec:validierung} verwendeten Metriken zusätzlich eine
manuell verfasste Referenzantwort ("Ground Truth"), gegen die Antwort und
Kontext inhaltlich abgeglichen werden.

\subsection{Die fünf RAGAS-Metriken}
\begin{itemize}
    \item \textbf{Faithfulness} -- "Ist die Antwort durch den Kontext gedeckt?"
    \begin{itemize}
        \item Braucht keine Referenz; zerlegt Antwort in atomare Claims, prüft jeden per NLI gegen den Kontext
        \item Niedriger Wert = Modell halluziniert
    \end{itemize}
    \item \textbf{AnswerRelevancy} -- "Beantwortet die Antwort die gestellte Frage?"
    \begin{itemize}
        \item Braucht keine Referenz; generiert Rückwärts-Fragen aus der Antwort, vergleicht Embedding-Ähnlichkeit zur Originalfrage
        \item Niedriger Wert = Antwort ausweichend/zu allgemein
    \end{itemize}
    \item \textbf{ContextPrecision} -- "Sind die abgerufenen Chunks relevant?"
    \begin{itemize}
        \item Braucht Referenzantwort; rangbasierte Average Precision über die Relevanz jedes abgerufenen Chunks
        \item Niedriger Wert = Retrieval liefert Rauschen
    \end{itemize}
    \item \textbf{ContextRecall} -- "Deckt der Kontext alles ab, was nötig wäre?"
    \begin{itemize}
        \item Braucht Referenzantwort; prüft, ob jede Aussage der Referenzantwort im Kontext belegbar ist
        \item Niedriger Wert = Retrieval verpasst wichtige Information
        \item Misst ausschließlich Retrieval-Qualität, nicht das antwortende LLM -- deshalb über verschiedene LLMs bei sonst gleicher Pipeline identisch (internes Konsistenz-Signal)
    \end{itemize}
    \item \textbf{FactualCorrectness} -- "Wie viele Fakten stimmen mit der Referenz überein?"
    \begin{itemize}
        \item Braucht Referenzantwort; F1-Score über atomare Fakten (generierte Antwort vs.\ Referenz)
        \item Durchgängig niedrigste Metrik (typisch 0,4--0,7) -- normal bei striktem Fakten-F1 gegen manuell formulierte Referenzen, nicht direkt mit den anderen vier Metriken vergleichbar
    \end{itemize}
    \item Alle Werte in $[0,1]$, höher ist besser; der in dieser Abgabe verwendete $\varnothing$-Score ist das ungewichtete Mittel aller fünf Metriken je Kombination (grobe Ranking-Kennzahl, für Detailanalysen sind Einzelmetriken aussagekräftiger)
\end{itemize}

\subsection{Evaluations-Setup}
\begin{itemize}
    \item Fester 12-Fragen-Datensatz mit Referenzantworten (\texttt{eval/dataset.py}), deckt alle vier Handbuch-Themenbereiche ab
    \item Gleicher Demo-Korpus und gleiches Fragenset über alle sechs Studien -- Vergleichbarkeit zwischen den Studien
    \item Die vier Pipeline-Stufen-Vergleiche (Parser, Chunking, Embedding, Retrieval) laufen mit fest auf \texttt{gpt-4o-mini} fixiertem Antwort-LLM und werden auch von \texttt{gpt-4o-mini} als Richter bewertet -- da das Antwort-LLM dort für alle verglichenen Varianten konstant bleibt, ist ein Self-Preference-Bias ausgeschlossen; ergänzend wurden die Ergebnisse für Retrieval und Embedding stichprobenartig gegen ein zweites, größeres Richtermodell (Mistral Large) validiert, ohne wesentliche Abweichungen zum primären Richter
    \item Der Antwort-LLM-Vergleich selbst (GPT-4o-mini, Mistral \texttt{mistral-small-latest}, GPT-4o) sowie das Best-of-Breed-Experiment testen gezielt mehrere Antwort-LLMs gegeneinander; da hier Modelle unterschiedlicher Anbieter direkt gegeneinander antreten, wird als Richter ein von allen Kandidaten unabhängiges Modell eingesetzt (Claude Sonnet 5, Anthropic)
    \item Gesamtumfang: 30 Kombinationen $\times$ 12 Fragen $\times$ 5 Metriken $=$ 1.800 Einzel-Scores
    \item Reproduzierbar über \texttt{scripts/evaluate\_*\_ragas.py}
\end{itemize}

\subsection{Bewertung der einzelnen Komponenten}

Mit der in den vorigen Abschnitten beschriebenen Methodik wurde jede der
fünf in Kapitel~\ref{sec:vergleiche} vorgestellten Dimensionen einzeln
ausgewertet -- jeweils mit allen übrigen Komponenten auf dem
Produktiv-Default fixiert. Die vollständigen Ergebnistabellen, Diagramme
und Rohdaten je Dimension sind nicht an dieser
Stelle ausgeführt, sondern im interaktiven Dashboard
\texttt{docs/dashboard.html} zusammengefasst. Ergänzend liegt zu jeder Dimension eine
eigene Dokumentationsdatei mit ausführlicher Diskussion vor:
\texttt{docs/PARSER.md}, \texttt{docs/CHUNKING.md},
\texttt{docs/EMBEDDING.md}, \texttt{docs/RETRIEVAL.md} und
\texttt{docs/LLM.md}. Im direkten Antwort-LLM-Vergleich unter dem unabhängigen Richter liegen GPT-4o-mini ($\varnothing$ 0,741) und GPT-4o ($\varnothing$ 0,739) praktisch gleichauf, mit Mistral ($\varnothing$ 0,662) deutlich dahinter. Abschließend wurde ein Best-of-Breed-Experiment durchgeführt, das die vier Einzelsieger kombiniert und prüft, ob sich deren Einzelverbesserungen tatsächlich addieren.

\subsection{Best-of-Breed-Experiment}
\label{sec:best-of-breed}

Das abschließende sechste Experiment kombiniert die vier Einzelsieger der
vorigen Vergleiche zum Zeitpunkt ihrer Ermittlung (Unstructured + Semantic
Chunking + \texttt{text-embedding-3-large} + Rerank) zu einer
Gesamt-Pipeline und testet sie auf dem Demo-Korpus über die drei
Antwort-LLMs (GPT-4o-mini, Mistral \texttt{mistral-small-latest}, GPT-4o) gegen die Produktiv-Baseline,
bewertet durch den in Abschnitt~\ref{sec:validierung} beschriebenen
unabhängigen Richter. Im Durchschnitt über alle drei Antwort-LLMs und
fünf Metriken schlägt Best-of-Breed die Baseline deutlich ($\varnothing$
0,819 vs.\ 0,716, $+$14,4\,\%). Der Vorsprung gilt jedoch nicht für jede
Einzelmetrik und jedes Antwort-LLM gleichermaßen: Mit GPT-4o gewinnt
Best-of-Breed alle fünf Metriken; mit Mistral sinkt die Faithfulness
gegenüber der Baseline (0,799 auf 0,746); mit GPT-4o-mini sinken sowohl
Faithfulness (0,906 auf 0,896) als auch FactualCorrectness (0,508 auf
0,479). Innerhalb von Best-of-Breed ist GPT-4o zugleich das mit Abstand
stärkste Antwort-LLM ($\varnothing$ 0,863, vor GPT-4o-mini mit 0,821 und
Mistral mit 0,772). Die methodische Kernaussage: Einzelverbesserungen
addieren sich nicht automatisch zu einer gleichmäßigen Verbesserung --
jede Kombination muss selbst gemessen werden. Details in
\texttt{docs/dashboard.html}, Tab "Best-of-Breed".

% =====================================================================
\section{Bekannte Grenzen, Ausblick und Fazit}
\label{sec:grenzen}
% =====================================================================

\subsection{Bekannte Grenzen und Ausblick}

\begin{itemize}
    \item Inkrementelles Ingestion-Update: Ein gezielter Fix ergänzt einen Datei-Fingerabdruck (SHA-256 je PDF, gespeichert in einem Manifest neben der Collection); seitdem verarbeitet jeder Ingestion-Lauf nur noch neue oder inhaltlich geänderte PDFs neu, entfernt gelöschte Dateien aus der Collection und überspringt unveränderte Dateien vollständig -- ein Lauf ohne Änderungen ist damit ein reiner No-op ohne API-Kosten
    \item Seitenbasiertes Header-Splitting: Ein gezielter Fix führt Abschnitte, die ohne neue Überschrift über eine Seitengrenze laufen, inzwischen mit dem vorigen Abschnitt zusammen (die Seitenangabe wird dann zu einem Bereich); der PDF-Deep-Link zeigt bei solchen zusammengeführten Chunks aber weiterhin nur auf die Startseite, nicht auf eine ggf.\ relevantere Folgeseite
    \item Speicherumfang bewusst begrenzt: aktuell 4 von 21 verfügbaren Abaqus-2017-Handbüchern eingelesen ($\sim$5.100 von $\sim$21.800 Seiten), um Ingestion-Dauer und Kosten für Entwicklung/Tests gering zu halten
    \item Vergleichsumfang der vier Einzeldimensionen auf den Demo-Korpus beschränkt: Die Parser-, Chunking-, Embedding- und Retrieval-Vergleiche laufen ausschließlich auf dem $\sim$53-Seiten-Demo-Korpus, nicht auf dem vollen Produktivkorpus -- wegen der Laufzeit von Unstructured (\texttt{hi\_res}) und Semantic Chunking auf $\sim$5.100 Seiten wurde eine Ausweitung im Rahmen dieser Abgabe bewusst nicht mehr durchgeführt. Das gilt auch für das Best-of-Breed-Experiment selbst, dessen Ergebnisse (Abschnitt~\ref{sec:best-of-breed}) sich ebenfalls auf den Demo-Korpus beziehen
\end{itemize}

\subsection{Fazit}
\label{sec:fazit}

Das im Rahmen dieser Arbeit entwickelte System zeigt, dass sich die
mehrere tausend Seiten umfassenden Abaqus-Handbücher erfolgreich zu einem
konversationsfähigen Chatbot verarbeiten lassen. Der Bot beantwortet
Bedienungsfragen ausschließlich auf Basis der Handbuchinhalte, macht jede
Antwort über eine seitengenaue Quellenangabe im Original nachprüfbar und
versteht dank eines einfachen Gesprächsgedächtnisses auch Folgefragen.

Über die reine Funktionsfähigkeit hinaus war ein zentrales Ziel der
Arbeit, die Designentscheidungen der Pipeline nicht nur zu begründen,
sondern auch empirisch zu überprüfen. Dazu wurde jede der fünf
austauschbaren Pipeline-Stufen -- Parser, Chunking, Embedding, Retrieval
und Antwort-LLM -- über eine eigene Registry konsequent austauschbar
gestaltet und in sechs unabhängigen RAGAS-Studien mit insgesamt 1.800
Einzel-Scores gegen Alternativen getestet. Damit liegen für alle fünf
Dimensionen belastbare Tendenzaussagen vor, statt sich allein auf
Intuition oder Einzelbeispiele zu verlassen.

Eine methodische Erkenntnis liefert dabei das abschließende
Best-of-Breed-Experiment: Die Kombination der vier jeweils besten
Einzelkomponenten schlägt die Produktiv-Baseline im Durchschnitt über alle
drei getesteten Antwort-LLMs und alle fünf Metriken deutlich
($+$14,4\,\%), jedoch nicht bei jeder Einzelmetrik gleichermaßen -- mit
Mistral (\texttt{mistral-small-latest}) und, geringfügig, auch mit GPT-4o-mini als Antwort-LLM sinkt die
Faithfulness gegenüber der Baseline leicht. Einzelverbesserungen addieren
sich also nicht automatisch zu einer gleichmäßigen Verbesserung, und selbst
ein neuer Gesamtsieger ersetzt nicht die Notwendigkeit, jede Kombination
einzeln zu messen. Für den praktischen Einsatz leitet sich daraus die
pragmatischste Empfehlung ab: mindestens Cross-Encoder-Reranking statt MMR
einsetzen -- der größte Einzelgewinn bei geringstem Umbauaufwand.
Innerhalb der Best-of-Breed-Pipeline erweist sich GPT-4o zudem als das mit
Abstand stärkste Antwort-LLM ($\varnothing$ 0,863 gegenüber 0,821 für
GPT-4o-mini und 0,772 für Mistral); die Produktiv-App bietet Best-of-Breed
dementsprechend nur auf dem vollen Produktivkorpus und standardmäßig mit
GPT-4o als Antwort-LLM an, trotz dessen höherer Kosten pro Token gegenüber
dem Produktiv-Default \texttt{gpt-4o-mini}.

Die in Abschnitt~\ref{sec:grenzen} genannte Grenze des auf den Demo-Korpus
beschränkten Vergleichsumfangs -- die auch für das Best-of-Breed-Experiment
selbst gilt -- schränkt die Übertragbarkeit der genauen Effektgrößen auf
den vollen Produktivkorpus ein; eine Validierung des Best-of-Breed-Experiments
auf dem vollen Korpus ist bewusst nicht Teil dieser Arbeit. Insgesamt zeigt die
Arbeit, dass ein systematischer, austauschbarer Pipeline-Aufbau in
Kombination mit einer konsistenten RAGAS-Evaluation es erlaubt,
RAG-Designentscheidungen empirisch statt nur intuitiv zu treffen.

\subsection{Hinweis zum Einsatz von KI-Werkzeugen}

Bei der Erstellung der vorliegenden Arbeit wurde KI-Unterstützung eingesetzt.

\end{document}
